\documentclass[11pt]{article}
\usepackage[T1]{fontenc}
\usepackage{amsmath,amssymb,amsthm}
\usepackage[margin=2.5cm]{geometry}

\newtheorem*{claim}{Claim (Subcase 2b)}
\newtheorem*{remark}{Remark}

\begin{document}

\section*{Proposition 4.1, Subcase 2b}

\paragraph{Setting.} The quadratic model is
\[
m_k(p) = f_k + g_k^\top p + \tfrac12 p^\top B_k p ,
\]
with $g_k \neq 0$, and the Cauchy point is $p_k^C = -\alpha_k^C g_k$, where
$\alpha_k^C$ minimizes $m_k(-\alpha g_k)$ over $0 \le \alpha \le \Delta_k/\|g_k\|$.
Define
\[
\varphi(\alpha) := m_k(-\alpha g_k) = f_k - \alpha \|g_k\|^2 + \tfrac12 \alpha^2\, g_k^\top B_k g_k .
\]
Case 2 is $g_k^\top B_k g_k > 0$; then $\varphi$ is a strictly convex parabola with
unconstrained minimizer
\[
\alpha^* = \frac{\|g_k\|^2}{g_k^\top B_k g_k} > 0 .
\]

\begin{claim}
If $g_k^\top B_k g_k > 0$ and $\alpha^* > \Delta_k/\|g_k\|$, then
$\alpha_k^C = \Delta_k/\|g_k\|$ and
\[
m_k(0) - m_k(p_k^C) \ \ge\ \tfrac12 \|g_k\|\,\Delta_k .
\]
\end{claim}

\begin{proof}
\textbf{Step 1: $\alpha_k^C = \Delta_k/\|g_k\|$.}
Since $\varphi'(\alpha) = -\|g_k\|^2 + \alpha\, g_k^\top B_k g_k$ and $g_k^\top B_k g_k>0$,
we have $\varphi'(\alpha) < 0$ for every $\alpha < \alpha^*$. Hence $\varphi$ is strictly
decreasing on $[0,\alpha^*]$. The feasible interval $[0, \Delta_k/\|g_k\|]$ is contained in
$[0,\alpha^*]$ (by hypothesis $\Delta_k/\|g_k\| < \alpha^*$), so the minimum of $\varphi$ over it
is attained at the right endpoint: $\alpha_k^C = \Delta_k/\|g_k\|$.

\textbf{Step 2: bound on the decrease.}
For any $\alpha \in [0,\alpha^*]$,
\[
\varphi(0) - \varphi(\alpha) = \alpha \|g_k\|^2 - \tfrac12 \alpha^2\, g_k^\top B_k g_k .
\]
Since $0 \le \alpha \le \alpha^*$ and $g_k^\top B_k g_k > 0$,
\[
\tfrac12 \alpha^2\, g_k^\top B_k g_k \ \le\ \tfrac12 \alpha\, \alpha^*\, g_k^\top B_k g_k
= \tfrac12 \alpha \|g_k\|^2 ,
\]
using $\alpha^* g_k^\top B_k g_k = \|g_k\|^2$. Therefore
\[
\varphi(0) - \varphi(\alpha) \ \ge\ \alpha\|g_k\|^2 - \tfrac12 \alpha\|g_k\|^2 = \tfrac12 \alpha \|g_k\|^2 .
\tag{$\ast$}
\]

\textbf{Step 3: conclusion.}
Applying $(\ast)$ with $\alpha = \alpha_k^C = \Delta_k/\|g_k\| \in [0,\alpha^*]$, and noting that
$\varphi(0)=m_k(0)$ and $\varphi(\alpha_k^C) = m_k(p_k^C)$,
\[
m_k(0) - m_k(p_k^C) \ \ge\ \tfrac12 \cdot \frac{\Delta_k}{\|g_k\|}\cdot \|g_k\|^2
= \tfrac12 \|g_k\|\,\Delta_k .
\]
Finally, since $\Delta_k \ge \min\!\left(\Delta_k, \|g_k\|/\|B_k\|\right)$, this yields
inequality (13) in this subcase.
\end{proof}

\begin{remark}[Meaning of ``the slope dominates'']
Inequality $(\ast)$ says that, as long as $\alpha$ does not exceed $\alpha^*$, the actual decrease is at
least half of the linear decrease predicted by the slope at the origin,
$-\varphi'(0)\,\alpha = \alpha\|g_k\|^2$. Indeed, the ratio of the curvature term
to the linear term is
\[
\frac{\tfrac12\alpha^2\, g_k^\top B_k g_k}{\alpha\|g_k\|^2} = \frac{\alpha}{2\alpha^*} \ \le\ \frac12
\qquad (0<\alpha\le\alpha^*).
\]
\end{remark}

\end{document}
